\section{Introduction}

Language contact commonly results in the borrowing of foreign lexical items
from a donor language. Borrowed words rarely enter unchanged; instead, they
undergo \emph{nativization}, adapting to the phonological, morphological, and
orthographic conventions of the recipient~\cite{haugen1950,wang2025}. The
phonological and morphological dimensions of loanword integration have been
studied extensively~\cite{kang2011, uffmann2015}, and the influence of
orthography on adaptation is also documented~\cite{wang2025,daland2015}.
However, \emph{orthographic nativization} has received considerably less
attention as a computational task, despite its relevance to text normalization,
spell-checking, and corpus construction in multilingual natural language
processing.

Within computational linguistics, loanword modeling efforts have targeted
related but distinct problems. Finite State Transducer (FST) cascades with Optimality-Theoretic
constraint ranking have modeled the morpho-phonological transformations of
Arabic loanwords entering Swahili for bilingual lexicon
generation~\cite{tsvetkov2015}. FSTs have also shown that much of the
phonological adaptation from English into Japanese can be captured by
finite-state mappings~\cite{mao2016}. More broadly, nativization in loanword
phonology has been shown to be systematic and rule-governed across
typologically diverse languages~\cite{kang2011}.

These studies focus on modeling phonological nativization patterns. What they
do not address is orthographic nativization as a \emph{generation} problem:
given a loanword, produce its standardized written form according to the
recipient language's orthographic conventions. This distinction matters, as
phonological and orthographic nativization are not
equivalent~\cite{wang2025,daland2015}. A phonological model produces an
abstract pronunciation representation; an orthographic nativization system must
produce a concrete character string that obeys the recipient's spelling rules.

We address this gap with a rule-based rewrite cascade drawn from prior FST work
on loanword modeling~\cite{mao2016}. By deriving rewrites from prescriptive
guidelines and linguistic literature, the approach is transparent and
linguistically grounded. It also sidesteps the limitation of data-driven
methods, which require parallel corpora typically unavailable for low-resource
languages, offering a practical middle ground.

We evaluate the proposed method using Filipino as a testbed. Filipino borrows
extensively from English across education, medicine, commerce, and public
administration~\cite{martin2020,baklanova2017}, and official nativization
guidelines are codified in~\cite{almario2014}. These properties enable
systematic evaluation against prescriptive orthographic standards. Filipino is
also typologically distinct from English in phonological inventory, syllable
structure, and orthographic depth~\cite{llamzon1966,schachter1972,llamzon1997},
ensuring the testbed exercises the full capacity of the proposed cascade.

This study makes the following contributions:
\begin{itemize}
      \item \textbf{A rewrite cascade framework for low-resource orthographic nativization}.
            We formalize orthographic nativization as a rewrite cascade,
            demonstrating that prescriptive linguistic rules can be
            operationalized as an interpretable computational pipeline
            without large parallel corpora.

      \item \textbf{An empirical case study on English--Filipino borrowing}.
            Using Filipino as a testbed, we implement and evaluate the framework on real
            loanword data, demonstrating practical effectiveness while identifying
            systematic sources of error.
\end{itemize}